% 50e Olympiade Internationale de Mathématiques — Brême, Allemagne.
% Énoncés officiels en français (imo-official.org). Chaque problème vaut 7 points.

\begin{center}
  \emph{Premier jour --- mercredi 15 juillet 2009 --- Durée : 4 heures 30}
\end{center}

\subsection*{Problème 1}
Soit $n$ un entier strictement positif et soit $a_1, \ldots, a_k$, avec $k \geqslant 2$, des entiers
strictement positifs distincts appartenant à l'ensemble $\{1, \ldots, n\}$ tels que $n$ divise
$a_i(a_{i+1} - 1)$ pour $i = 1, \ldots, k - 1$.

Montrer que $n$ ne divise pas $a_k(a_1 - 1)$.

\subsection*{Problème 2}
Soit $ABC$ un triangle et $O$ le centre de son cercle circonscrit. Les points $P$ et $Q$ sont des points
intérieurs aux côtés $CA$ et $AB$ respectivement. Soit $K$, $L$ et $M$ les milieux respectifs des segments
$BP$, $CQ$ et $PQ$, et soit $\Gamma$ le cercle passant par $K$, $L$ et $M$. On suppose que la droite $(PQ)$
est tangente au cercle $\Gamma$.

Montrer que $OP = OQ$.

\subsection*{Problème 3}
Soit $s_1, s_2, s_3, \ldots$ une suite strictement croissante d'entiers strictement positifs telle que les
sous-suites
\[
s_{s_1}, s_{s_2}, s_{s_3}, \ldots \qquad \text{et} \qquad s_{s_1+1}, s_{s_2+1}, s_{s_3+1}, \ldots
\]
soient deux progressions arithmétiques.

Montrer que la suite $s_1, s_2, s_3, \ldots$ est aussi une progression arithmétique.

\bigskip
\begin{center}
  \emph{Second jour --- jeudi 16 juillet 2009 --- Durée : 4 heures 30}
\end{center}

\subsection*{Problème 4}
Soit $ABC$ un triangle tel que $AB = AC$. Les bissectrices de $\widehat{CAB}$ et $\widehat{ABC}$ rencontrent
respectivement les côtés $BC$ et $CA$ en $D$ et $E$. Soit $K$ le centre du cercle inscrit dans le triangle
$ADC$. On suppose que $\widehat{BEK} = 45^\circ$.

Trouver toutes les valeurs possibles de $\widehat{CAB}$.

\subsection*{Problème 5}
Déterminer toutes les fonctions $f$ de l'ensemble des entiers strictement positifs dans l'ensemble des
entiers strictement positifs telles que, pour tous entiers strictement positifs $a$ et $b$, il existe un
triangle non aplati dont les longueurs des côtés sont
\[
a, \ f(b) \ \text{ et } \ f(b + f(a) - 1).
\]

\subsection*{Problème 6}
Soit $a_1, a_2, \ldots, a_n$ des entiers strictement positifs distincts et soit $M$ un ensemble de $n - 1$
entiers strictement positifs ne contenant pas $s = a_1 + a_2 + \cdots + a_n$. Une sauterelle doit faire des
sauts le long de l'axe réel ; partant du point $0$, elle doit effectuer $n$ sauts vers la droite de longueurs
$a_1, a_2, \ldots, a_n$ dans l'ordre de son choix.

Montrer que la sauterelle peut choisir l'ordre de ses sauts de façon à ne passer par aucun point de $M$.
